\documentclass[11pt]{article}

%------------------------- fonts & encoding -----------------------------------
\usepackage[T1]{fontenc}
\usepackage[utf8]{inputenc}
\usepackage{amsmath,amssymb,amsthm}
\usepackage{mathtools}
\usepackage[sc]{mathpazo}
\linespread{1.06}
\usepackage{microtype}

%------------------------- layout ---------------------------------------------
\usepackage[a4paper,margin=1.05in,marginparwidth=1.0in]{geometry}
\usepackage{booktabs}
\usepackage{enumitem}
\setlist{itemsep=2pt,topsep=3pt}
\usepackage{xcolor}
\usepackage{graphicx}
\usepackage{tikz}
\usepackage{pgfplots}
\pgfplotsset{compat=1.17}
\usepackage[most,breakable]{tcolorbox}

%------------------------- palette --------------------------------------------
\definecolor{crimson}{HTML}{A51C30}
\definecolor{ink}{HTML}{1A1A1A}
\definecolor{steel}{HTML}{2C5F8A}
\definecolor{amber}{HTML}{9A6A00}
\definecolor{teal}{HTML}{1F6F6B}
\definecolor{paperlt}{HTML}{F7F4EF}

%------------------------- section styling ------------------------------------
\usepackage{titlesec}
\titleformat{\section}{\Large\bfseries\color{steel}}{\thesection}{0.7em}{}
\titleformat{\subsection}{\large\bfseries\color{ink}}{\thesubsection}{0.6em}{}
\titlespacing*{\section}{0pt}{1.6em}{0.6em}

%------------------------- hyperref -------------------------------------------
\usepackage[colorlinks=true,linkcolor=steel,citecolor=steel,urlcolor=crimson]{hyperref}

%------------------------- theorem environments -------------------------------
\theoremstyle{definition}
\newtheorem{definition}{Definition}[section]
\newtheorem{example}{Example}[section]
\newtheorem{exercise}{Exercise}[section]
\theoremstyle{plain}
\newtheorem{proposition}{Proposition}[section]
\theoremstyle{remark}
\newtheorem*{remark}{Remark}

%------------------------- callout boxes --------------------------------------
\newtcolorbox{intuition}{
  enhanced, breakable, colback=crimson!5, colframe=crimson, boxrule=0.5pt,
  left=8pt,right=8pt,top=6pt,bottom=6pt, arc=2pt,
  colbacktitle=crimson, coltitle=paperlt, fonttitle=\bfseries, title={Intuition}}

\newtcolorbox{pitfall}{
  enhanced, breakable, colback=amber!7, colframe=amber, boxrule=0.5pt,
  left=8pt,right=8pt,top=6pt,bottom=6pt, arc=2pt,
  colbacktitle=amber, coltitle=paperlt, fonttitle=\bfseries, title={Pitfall}}

\newtcolorbox{keyresult}{
  enhanced, breakable, colback=steel!6, colframe=steel, boxrule=0.5pt,
  left=8pt,right=8pt,top=6pt,bottom=6pt, arc=2pt,
  colbacktitle=steel, coltitle=paperlt, fonttitle=\bfseries, title={Takeaway}}

\newtcolorbox{observed}{
  enhanced, breakable, colback=teal!6, colframe=teal, boxrule=0.5pt,
  left=8pt,right=8pt,top=6pt,bottom=6pt, arc=2pt,
  colbacktitle=teal, coltitle=paperlt, fonttitle=\bfseries, title={Measured}}

%------------------------- math operators -------------------------------------
\DeclareMathOperator{\Var}{Var}
\DeclareMathOperator{\Cov}{Cov}
\DeclareMathOperator{\Corr}{corr}
\DeclareMathOperator{\diag}{diag}
\DeclareMathOperator{\std}{std}
\DeclareMathOperator*{\avg}{\mathbb{E}}
\newcommand{\T}{^{\!\top}}
\newcommand{\R}{\mathbb{R}}
\newcommand{\half}{\tfrac12}
\newcommand{\wvec}{\mathbf{w}}
\newcommand{\Xmat}{\mathbf{X}}
\newcommand{\Fmat}{\mathbf{F}}
\newcommand{\Vmat}{\mathbf{V}}
\newcommand{\Umat}{\mathbf{U}}
\newcommand{\Rmat}{\mathbf{R}}
\newcommand{\Lam}{\boldsymbol{\Lambda}}
\newcommand{\Del}{\boldsymbol{\Delta}}
\newcommand{\fvec}{\mathbf{f}}

%------------------------- title ----------------------------------------------
\title{\vspace{-2.2em}
  {\color{steel}\rule{\textwidth}{1.3pt}}\\[0.7em]
  \Huge\bfseries The Factor Covariance Forecast\\[0.35em]
  \large\mdseries Turning a history of factor returns into a forward-looking risk matrix --- in four layers\\[0.4em]
  {\color{steel}\rule{\textwidth}{0.7pt}}}
\author{}
\date{}

%------------------------- watermark footer -----------------------------------
\usepackage{fancyhdr}
\pagestyle{fancy}
\fancyhf{}
\renewcommand{\headrulewidth}{0pt}
\renewcommand{\footrulewidth}{0pt}
\fancyfoot[L]{\footnotesize\color{black!45}MSCI Barra USE4 Imitation}
\fancyfoot[R]{\footnotesize\color{black!45}\thepage}

\begin{document}
\maketitle
\thispagestyle{empty}

\begin{abstract}
The cross-sectional regression hands us a history of factor returns; this stage turns that history into a single forward-looking object: the factor covariance matrix $\Fmat$, the $68\times68$ array of variances and covariances that drives all factor risk in $\sigma_p^2 = \wvec\T\Xmat\Fmat\Xmat\T\wvec + \wvec\T\Del\wvec$. Estimating a covariance sounds routine -- take the sample covariance of the returns and stop -- but the sample covariance is a poor \emph{forecast}, and the ways it fails are specific and correctable. It weights a five-year-old crisis the same as last week; it ignores the serial correlation that makes a month riskier than twenty-one independent days; its eigenvalues are systematically too spread out, so the portfolios it thinks are safest are the ones it most under-forecasts; and its overall level lags regime changes. USE4 addresses each failure with a dedicated layer: an exponentially-weighted base with separate memories for volatilities and correlations, a Newey--West correction for serial correlation, a Monte-Carlo eigenfactor adjustment that de-biases the eigenvalues, and a volatility-regime scaling that tracks the current level. This note explains the mathematics of each layer, the statistical reason it is needed, and -- because the matrix exists to be optimised against -- why each one matters in practice. The eigenfactor adjustment is the headline: it is the single most important correction for a covariance matrix that will be fed to an optimiser.
\end{abstract}

\begin{tcolorbox}[colback=paperlt,colframe=ink!30,boxrule=0.4pt,arc=2pt,left=10pt,right=10pt,top=8pt,bottom=8pt]
\textbf{How to read these notes.} This is a long set, because the covariance forecast is where the model's statistics are deepest. We build it in four layers (Sections~\ref{sec:ewma}--\ref{sec:vra}), each fixing one concrete failure of the naive sample covariance, then check that it works (Section~\ref{sec:valid}). Read Section~\ref{sec:eigen}, the eigenfactor adjustment, slowest; it carries the weight of the set and is the correction that most changes how the matrix behaves in an optimiser. The numbers labelled as measured come from a single run of our personal build.

\smallskip
Four coloured boxes recur: \textcolor{crimson}{\textbf{Intuition}} gives the mental picture, \textcolor{amber}{\textbf{Pitfall}} flags what breaks without the fix, \textcolor{steel}{\textbf{Takeaway}} states a load-bearing result, and \textcolor{teal}{\textbf{Measured}} reports real numbers from the build run.

\smallskip
\textbf{Reference material.}
\begin{itemize}[leftmargin=1.4em,topsep=2pt]
  \item External: Menchero, Orr \& Wang (2011), USE4 Methodology Notes (Sec. 4); Menchero, Wang \& Orr (2011), \textit{Eigen-adjusted covariance matrices}.
\end{itemize}
\end{tcolorbox}

\tableofcontents
\bigskip

\section{What the covariance is, and why it is the hard part}\label{sec:what}

\paragraph{Learning objectives.}
\begin{itemize}[leftmargin=1.4em,topsep=2pt]
  \item State what $\Fmat$ is and where it enters the risk model.
  \item Explain why a \emph{forecast} of covariance is harder than a \emph{measurement} of it.
\end{itemize}

The factor covariance matrix $\Fmat$ is the engine of every risk number the model produces. A portfolio's factor risk is $\wvec\T\Xmat\Fmat\Xmat\T\wvec$: the portfolio weights $\wvec$ map through the exposures $\Xmat$ into factor space, and $\Fmat$ -- the $68\times68$ matrix of factor variances and covariances -- converts those factor positions into a variance. Get $\Fmat$ wrong and every forecast volatility, every tracking error, every optimised portfolio is wrong with it. It is the most consequential single object in the model.

It is also the hardest to get right, and the reason is the gap between \emph{measuring} and \emph{forecasting}. Measuring a covariance is easy: take the realised factor returns and compute their sample covariance. But the sample covariance describes the past, and risk management needs the future. The sample covariance is a biased and noisy forecast of next month's covariance, and -- crucially -- it is biased in \emph{structured}, well-understood ways. It treats every day in its window as equally informative, so it is slow to react when volatility jumps. It assumes the daily returns are independent, so it misses the serial correlation that makes a calendar month riskier than its days suggest. Its eigenvalues are too dispersed, an artefact of estimation error that bites hardest exactly where an optimiser is most sensitive. And its level drifts behind the current regime. Each of these is a known failure with a known fix, and the four layers of this stage are those four fixes, applied in order.

\begin{intuition}
A sample covariance is a rear-view mirror: a faithful record of where the road has been. To drive you need a windscreen. The four layers reshape the mirror's image into a forecast -- weighting recent road more, accounting for the bends that one frame misses, correcting a lens distortion that makes some directions look flatter than they are, and adjusting for how fast the scenery is currently changing. The raw image is not wrong; it is just not yet a forecast.
\end{intuition}

\begin{keyresult}
$\Fmat$ is the $68\times68$ factor covariance that drives all factor risk. The sample covariance is a poor forecast of it in four specific, correctable ways; the build fixes each with a dedicated layer. Everything in this note is one of those four corrections.
\end{keyresult}

\section*{Notation}
\addcontentsline{toc}{section}{Notation}

\begin{center}
\begin{tabular}{@{}ll@{}}
\toprule
Symbol & Meaning \\
\midrule
$\Fmat$ & the $68\times68$ factor covariance forecast (monthly horizon) \\
$f_{k,s}$ & return of factor $k$ on day $s$ (from the daily cross-sectional regression) \\
$\sigma_k,\ R_{ij}$ & factor volatility and factor--factor correlation \\
$h$ & a half-life: the lag at which an exponential weight falls to one-half \\
$w_s$ & EWMA weight on day $s$, $\propto \exp(-\ln 2\,(t-s)/h)$, normalised \\
$\Vmat_\ell$ & lag-$\ell$ autocovariance matrix of the daily factor returns \\
$L$ & number of Newey--West lags \\
$\lambda_k,\ \mathbf{u}_k$ & the $k$-th eigenvalue and eigenvector (eigenfactor) of the covariance \\
$b_k$ & bias statistic of factor (or eigenfactor) $k$: realised risk $/$ predicted risk \\
$B_t,\ \lambda_{\text{vol}}$ & daily cross-sectional bias, and the volatility-regime multiplier \\
$\sigma_p^2$ & portfolio variance $=\wvec\T\Xmat\Fmat\Xmat\T\wvec+\wvec\T\Del\wvec$ \\
\bottomrule
\end{tabular}
\end{center}

\section{Why daily, and the four-layer plan}\label{sec:plan}

\paragraph{Learning objectives.}
\begin{itemize}[leftmargin=1.4em,topsep=2pt]
  \item Explain why the covariance is estimated from daily, not monthly, factor returns.
  \item Lay out the four layers and what each one corrects.
\end{itemize}

The covariance is built from \emph{daily} factor returns, not the monthly ones used elsewhere. The reason is sample size. The monthly cross-sectional regression yields a few hundred monthly return observations -- too few to estimate a $68\times68$ matrix well, and far too few for the later corrections, which lean on having many observations: the serial-correlation correction reads daily autocorrelation, and the eigenvalue de-biasing needs enough data for its simulations to mean something. So the model runs the \emph{same} constrained, $\sqrt{\text{mcap}}$-weighted regression once per trading day, holding the month-end exposures fixed within each month, producing thousands of daily factor-return observations. The daily series reconciles to the monthly one -- summing the daily factor returns across a month reproduces the monthly factor return very closely -- so nothing is lost in the move to daily; a great deal of statistical resolution is gained.

From the daily returns the forecast is assembled in four layers, each correcting one failure of the naive estimate:
\begin{enumerate}[leftmargin=1.6em,topsep=2pt]
  \item \textbf{EWMA base} -- weight recent days more than old ones, with separate memories for volatilities and correlations (Section~\ref{sec:ewma}).
  \item \textbf{Newey--West} -- add back the serial correlation that a day-by-day variance misses, then scale to the monthly horizon (Section~\ref{sec:nw}).
  \item \textbf{Eigenfactor adjustment} -- de-bias the eigenvalues, which the sample spreads too wide (Section~\ref{sec:eigen}).
  \item \textbf{Volatility-regime scaling} -- lift or lower the whole matrix to match the current level of market volatility (Section~\ref{sec:vra}).
\end{enumerate}
The order matters: each layer takes the previous layer's matrix as its input, and the eigenvalue de-biasing in particular must act on a matrix that already carries the serial-correlation correction.

\begin{pitfall}
Estimating a $68\times68$ covariance from a few hundred monthly returns is the classic small-sample trap: the matrix has $\tfrac{68\cdot 69}{2}=2{,}346$ free entries, and a few hundred observations cannot pin them down. The estimate would be dominated by noise and might not even be usable in an optimiser. Moving to daily data lifts the ratio of observations to parameters by an order of magnitude -- the precondition for everything that follows.
\end{pitfall}

\section{Layer 1: the EWMA base}\label{sec:ewma}

\paragraph{Learning objectives.}
\begin{itemize}[leftmargin=1.4em,topsep=2pt]
  \item Write the exponentially-weighted covariance and explain the half-life.
  \item Explain why volatilities and correlations are given \emph{different} half-lives.
\end{itemize}

The base estimate is an exponentially-weighted moving average (EWMA) of the daily factor returns' co-movements. Where an equal-weighted sample covariance averages $(f_{i,s}-\bar f_i)(f_{j,s}-\bar f_j)$ with the same weight on every day $s$, the EWMA discounts the past geometrically:
\[
\Fmat^{0}_{ij} \;=\; \sum_{s\le t} w_s\,(f_{i,s}-\bar f_i)(f_{j,s}-\bar f_j),
\qquad
w_s \;\propto\; \exp\!\Big(-\ln 2\,\tfrac{t-s}{h}\Big),
\]
with the weights normalised to sum to one. The half-life $h$ is the knob: a day exactly $h$ days in the past carries half the weight of today, a day $2h$ back a quarter, and so on. A short half-life makes the estimate nimble but jumpy; a long one makes it smooth but slow. This is the same exponential-memory idea the Beta and Momentum factors used on single series, now applied to the whole matrix of factor co-movements.

The refinement that matters is using \emph{two} half-lives. Volatilities and correlations do not move at the same speed. Volatility is famously fast and clustering -- it spikes within days when a crisis hits and subsides gradually -- whereas the correlation \emph{structure} among factors is far more stable, drifting over quarters rather than lurching in a week. So the build estimates the two with different memories: a \emph{short} half-life for the volatilities $\sigma_i$ (a faster window, so the diagonal reacts quickly to a volatility spike) and a \emph{long} half-life for the correlations $R_{ij}$ (a slower window, so the off-diagonal structure stays steady and is not whipped around by noise). The two are then recombined into a covariance,
\[
\Fmat^{0}_{ij} \;=\; R_{ij}\,\sigma_i\,\sigma_j ,
\]
taking the volatilities from the fast estimate and the correlations from the slow one. The result reacts promptly in \emph{level} while holding a stable \emph{shape}.

\begin{intuition}
Think of volatility as the weather and correlation as the climate. The weather can change in a day -- you want a thermometer that responds fast. The climate shifts over seasons -- you want a long average so a single hot afternoon does not convince you the climate has changed. Splitting the half-lives gives each the memory that suits it: a fast thermometer for volatilities, a slow climatology for correlations.
\end{intuition}

\begin{keyresult}
The base is an EWMA covariance: recent days weighted more, with a short half-life for volatilities (which move fast) and a long half-life for correlations (which move slowly), recombined as $\Fmat^0_{ij}=R_{ij}\sigma_i\sigma_j$.
\end{keyresult}

\section{Layer 2: Newey--West and serial correlation}\label{sec:nw}

\paragraph{Learning objectives.}
\begin{itemize}[leftmargin=1.4em,topsep=2pt]
  \item Explain why a monthly variance is not just twenty-one daily variances.
  \item Write the Newey--West estimator and say why its weights keep the matrix valid.
\end{itemize}

The base covariance is a \emph{daily} quantity, but risk is forecast at a monthly horizon, and here a subtle error lurks. If daily returns were independent, a month's variance would be just the sum of its days' variances -- scale the daily covariance by the number of trading days ($\approx 21$) and stop. But daily factor returns are \emph{not} independent. Microstructure effects, stale prices in less-liquid names, and lead--lag relationships leave the daily series serially correlated, and that autocorrelation changes the monthly variance. For a monthly return that is the sum of daily returns $S=\sum_{d=1}^{n} f_d$,
\[
\Var(S) \;=\; \sum_{d}\Var(f_d) \;+\; 2\sum_{\ell\ge 1}\sum_d \Cov(f_d, f_{d-\ell}),
\]
so positive autocorrelation (the common case) makes the true monthly variance \emph{larger} than $n$ times the daily variance. Scaling the daily covariance by $21$ and ignoring the cross terms would under-forecast monthly risk.

The Newey--West estimator restores the missing cross terms. It adds the lagged autocovariance matrices $\Vmat_\ell$ -- the covariance of today's factor returns with those $\ell$ days earlier -- back into the estimate, but with declining (Bartlett) weights and only out to a modest number of lags $L$:
\[
\Fmat^{\text{NW}} \;=\; \Vmat_0 \;+\; \sum_{\ell=1}^{L}\Big(1-\tfrac{\ell}{L+1}\Big)\big(\Vmat_\ell + \Vmat_\ell\T\big),
\]
after which the whole matrix is scaled from a daily to a monthly horizon. The triangular weights $1-\ell/(L+1)$ are not cosmetic: they are exactly what guarantees the corrected matrix stays positive semi-definite -- a non-negotiable property, since a covariance with a negative eigenvalue would imply a portfolio with negative variance and would break any optimiser. A small number of lags suffices; the autocorrelation in daily factor returns dies out quickly, and adding distant, noisy lags would cost more in variance than it returns in bias.

\begin{intuition}
A month is not twenty-one unrelated coin flips; it is twenty-one days that lean on each other. If a factor tends to keep moving the same way for a few days, the month's swings are bigger than independent days would produce -- a fact a day-by-day variance cannot see. Newey--West measures how each day echoes the days just before it and adds those echoes back, with nearer echoes counting more, so the monthly forecast captures the persistence the daily view misses.
\end{intuition}

\begin{pitfall}
Annualising or ``monthlyising'' a daily covariance by a flat $\times 21$ is the trap. It silently assumes zero autocorrelation, and for positively autocorrelated factor returns it under-forecasts monthly risk -- the most dangerous direction to be wrong. The Bartlett weights matter too: drop them for flat weights and the autocovariance correction can push the matrix non-PSD, at which point it is not a covariance at all.
\end{pitfall}

\section{Layer 3: the eigenfactor adjustment}\label{sec:eigen}

\paragraph{Learning objectives.}
\begin{itemize}[leftmargin=1.4em,topsep=2pt]
  \item State the eigenvalue-dispersion bias of a sample covariance and why it is structured, not random.
  \item Explain the Monte-Carlo de-biasing and read the bias smile (Figure~\ref{fig:fcov_smile}).
  \item Say why this is the single most important correction for an optimised model.
\end{itemize}

This is the layer that matters most, and the one with the most surprising statistics. Diagonalise the covariance into its eigenvalues and eigenvectors. The eigenvectors are \emph{eigenfactors} -- uncorrelated portfolios of the original factors -- and each eigenvalue is the variance the matrix assigns to its eigenfactor. The trouble is that the sample's eigenvalues are systematically \emph{too dispersed}: estimation error pushes the largest sample eigenvalues above their true values and the smallest ones below. This is not bad luck on a given month; it is a structural property of estimating a covariance from finite data, and it always leans the same way.

The consequence is precise and damaging. The eigenfactors with the smallest eigenvalues are the portfolios the matrix believes are \emph{least} risky -- and those are exactly the ones whose risk it \emph{under}-estimates. We can see it directly by backtesting: form each eigenfactor from a forecast, then measure its realised volatility over the following period and compare to the predicted $\sqrt{\lambda_k}$. The ratio, the \emph{bias statistic} $b_k=\text{realised}/\text{predicted}$, should be one if the forecast were unbiased. Before correction it is not (Figure~\ref{fig:fcov_smile}, left): the lowest-variance eigenfactors realise up to nearly twice the risk the matrix predicted, while the highest-variance ones are mildly over-predicted -- a pronounced downward ``smile'' in $b_k$ across the eigenfactor spectrum.

\begin{intuition}
Why are the smallest eigenvalues too small? Estimation noise always finds a way to make some combination of factors look quieter than it truly is -- by chance, some eigenfactor's past wobbles partly cancelled. The fit rewards that accident with a tiny eigenvalue, but the cancellation will not repeat out of sample, so the eigenfactor's real risk is higher than its tiny predicted risk. The same noise inflates the largest eigenvalue. The sample spectrum is a funhouse mirror: it stretches the extremes, making the calmest directions look calmer and the wildest wilder than they are.
\end{intuition}

This matters enormously because $\Fmat$ exists to be \emph{optimised against}. A mean-variance optimiser hunts for low-risk ways to hold a view, so it is drawn precisely to the low-variance eigenfactors -- the very directions the matrix under-forecasts. Left uncorrected, the optimiser systematically loads into portfolios that look safe and are not, and the realised risk of an ``optimised'' book runs above its forecast. Correcting the smallest eigenvalues is therefore not a refinement; it is the difference between a covariance that is safe to optimise against and one that quietly misleads.

The correction is a Monte-Carlo eigenvalue de-biasing. Its logic is a bootstrap: take the estimated covariance and \emph{treat it as truth}; simulate many independent return panels drawn from it; re-estimate a covariance on each simulated panel; and, because the truth is now known, measure exactly how the estimation distorts each eigenvalue -- the average ratio of estimated to true eigenvalue, eigenfactor by eigenfactor. That ratio is the bias-correction profile: below one for the smallest eigenvalues (they come out too small), above one for the largest. The real sample's eigenvalues are then rescaled by the inverse of this profile -- inflating the small ones, gently deflating the large -- with a modest extra scaling that tunes the strength of the correction, and the matrix is reassembled from the de-biased eigenvalues and the original eigenvectors. The effect is visible in Figure~\ref{fig:fcov_smile}: after de-biasing, the smile flattens onto one, the previously badly under-forecast low-variance eigenfactors landing back inside the confidence band.

\begin{figure}[htbp]
\centering
\includegraphics[width=\linewidth]{figures/fcov_bias_stats.png}
\caption{The eigenfactor bias, before and after de-biasing, over the monthly backtest. \emph{Left:} the bias statistic $b_k=\text{realised}/\text{predicted risk}$ across the eigenfactors, ordered from lowest to highest variance. Before the adjustment the lowest-variance eigenfactors realise up to $\approx 1.95\times$ their predicted risk while the highest-variance ones are mildly over-predicted ($b<1$) -- the downward ``smile'' of a too-dispersed spectrum -- and after the Monte-Carlo de-biasing the whole curve flattens into the $\pm 2\!\cdot\!\text{CI}$ band around one. \emph{Right:} the per-factor bias for the country and style factors, all close to one (mildly conservative). The mean absolute deviation $|b-1|$ falls from $\approx 0.31$ to $\approx 0.07$.}
\label{fig:fcov_smile}
\end{figure}

\begin{keyresult}
A sample covariance's eigenvalues are too dispersed -- smallest too small, largest too large -- so it under-forecasts the risk of its lowest-variance eigenfactors, exactly the portfolios an optimiser favours. The Monte-Carlo eigenfactor adjustment measures the bias by simulating from the estimate and rescales the eigenvalues to remove it, flattening the bias smile. It is the single most important correction for a covariance used in optimisation.
\end{keyresult}

\section{Layer 4: the volatility-regime adjustment}\label{sec:vra}

\paragraph{Learning objectives.}
\begin{itemize}[leftmargin=1.4em,topsep=2pt]
  \item Explain why even an EWMA can lag a regime change in overall level.
  \item Write the volatility-regime multiplier and describe what it does.
\end{itemize}

The first three layers fix the \emph{shape} and the structural biases of the matrix, but its overall \emph{level} can still lag. An EWMA reacts to a volatility jump only as fast as its half-life allows, and at the onset of a crisis even a fast half-life trails the truth by days -- precisely when an accurate level matters most. The volatility-regime adjustment is a final, fast-reacting scalar that nudges the whole matrix up or down to match the current environment.

It is driven by the model's own recent forecasting error, aggregated across factors. On each day $t$ form the cross-sectional bias
\[
B_t \;=\; \frac{1}{68}\sum_{k=1}^{68}\Big(\frac{f_{k,t}}{\sigma_{k,t}}\Big)^{2},
\]
the average squared \emph{standardised} factor return. If the factors are, on average, moving exactly as much as the matrix predicted, each $f_{k,t}/\sigma_{k,t}$ is about one in magnitude and $B_t\approx 1$. If the factors are collectively moving more than predicted -- the signature of a volatility regime the matrix has not yet caught up to -- then $B_t>1$. A causal exponential average of $B_t$ (a short half-life, so it tracks the regime promptly) gives the current multiplier, and the matrix is scaled by it:
\[
\lambda_{\text{vol}} \;=\; \sqrt{\operatorname{EWMA}(B_t)},
\qquad
\Fmat \;\leftarrow\; \lambda_{\text{vol}}^{2}\,\Fmat .
\]
When recent factor moves have outrun the forecast, $\lambda_{\text{vol}}>1$ and the whole matrix is lifted; in calm periods $\lambda_{\text{vol}}<1$ and it is trimmed. Because it multiplies the entire matrix by a scalar, it changes the level without touching the carefully-built correlation shape or the de-biased eigenstructure. The result is a forecast that tracks volatility regimes promptly (Figure~\ref{fig:fcov_vol}): the predicted market volatility climbs steeply in 2008 and 2020 and recedes in quiet years.

\begin{figure}[htbp]
\centering
\includegraphics[width=\linewidth]{figures/fcov_vol.png}
\caption{The covariance forecast in action: the predicted monthly volatility of the country (market) factor, $\sqrt{\Fmat_{cc}}$, over the backtest. It sits near a median of $\approx 4\%$ per month ($\approx 14\%$ annualised) in normal times and spikes toward $\approx 24\%$ in crises (2008, 2020), driven up by the fast EWMA volatilities and the volatility-regime multiplier together. A risk forecast that did not move like this would badly mis-state risk exactly when it matters.}
\label{fig:fcov_vol}
\end{figure}

\begin{intuition}
The volatility-regime multiplier is the model grading its own recent homework. It asks: over the last few weeks, did the factors move as much as I said they would? If they moved more, the multiplier concludes the world has gotten riskier than the slow averages have registered and scales the whole forecast up; if less, it scales down. It is a fast corrective layered on top of the slower memory beneath -- reactive where the EWMA is deliberate.
\end{intuition}

\section{Does it work? bias statistics and PSD}\label{sec:valid}

\paragraph{Learning objectives.}
\begin{itemize}[leftmargin=1.4em,topsep=2pt]
  \item Define the bias statistic and state the calibration achieved.
  \item Confirm the forecast is always a valid covariance.
\end{itemize}

The whole pipeline is validated the same way the eigenfactor smile was diagnosed -- with \emph{bias statistics} on a long out-of-sample backtest. For each factor (or eigenfactor) the bias statistic compares realised to predicted risk over the history; a well-calibrated forecast scores near one, a forecast that under-predicts scores above one, and one that over-predicts scores below. Because nothing in the pipeline is fitted to the outcomes -- the half-lives, the lag count, and the eigenfactor scaling are fixed priors, and every forecast uses only data available at the time -- the backtest is genuinely out-of-sample, and splitting it into disjoint early and late eras gives essentially the same calibration, confirming the result is structural rather than a full-sample artefact.

\begin{observed}
On the build run the forecast is assembled from $\approx 6{,}300$ daily factor-return days into $288$ monthly forecasts (after a one-year warm-up), each a $68\times68$ symmetric, positive-semi-definite matrix. The per-factor bias statistic averages $\approx 0.95$ (median $\approx 0.97$, country $\approx 0.91$) -- inside the $[0.85,1.15]$ target band and mildly conservative. The eigenfactor de-biasing is the headline: the mean absolute bias deviation across eigenfactors falls from $\approx 0.31$ to $\approx 0.07$ (a roughly four-fold flattening), the ten lowest-variance eigenfactors from $\approx 0.61$ to $\approx 0.11$, with $\approx 97\%$ of eigenfactors landing within two confidence intervals of one. The volatility-regime multiplier averages $\approx 1.02$ but ranges from $\approx 0.69$ in calm to $\approx 2.2$ in crisis. The predicted country (market) volatility runs near $\approx 4\%$ monthly in normal times and spikes toward $\approx 24\%$ in crises; the latest forecast is $\approx 4.7\%$ monthly ($\approx 16\%$ annualised). Every forecast is PSD, with strictly positive eigenvalues throughout.
\end{observed}

\begin{exercise}
A factor's daily returns have positive first-order autocorrelation. You forecast its monthly variance as $21$ times its daily variance. Is your forecast too high or too low, and which term in $\Var\!\big(\sum_d f_d\big)$ did you omit?
\end{exercise}

\begin{exercise}
An optimiser is handed a covariance whose smallest eigenfactor has predicted volatility well below its true volatility. Explain, in two sentences, why the optimiser is especially likely to take a large position aligned with that eigenfactor, and what happens to the realised risk of the resulting portfolio.
\end{exercise}

\begin{exercise}
The volatility-regime bias is $B_t=\frac1{68}\sum_k (f_{k,t}/\sigma_{k,t})^2$. Suppose on a given day every factor moved exactly twice its predicted volatility. Compute $B_t$ and the resulting one-day multiplier $\lambda_{\text{vol}}=\sqrt{B_t}$, and say what the matrix would be scaled by.
\end{exercise}

\section{Summary, and one caveat}\label{sec:summary}

\paragraph{Learning objectives.}
\begin{itemize}[leftmargin=1.4em,topsep=2pt]
  \item Recite the four layers and what each corrects.
  \item State the central limitation of the forecast.
\end{itemize}

The build, end to end: run the cross-sectional regression daily to get a long history of daily factor returns; form an EWMA base covariance with a fast half-life for volatilities and a slow one for correlations; apply a Bartlett-weighted Newey--West correction to restore the serial correlation a daily variance misses, and scale to the monthly horizon; de-bias the eigenvalues with a Monte-Carlo eigenfactor adjustment, inflating the too-small eigenvalues that an optimiser would otherwise exploit; and finish with a volatility-regime multiplier that tracks the current level. The output is a positive-semi-definite, well-calibrated, regime-aware $68\times68$ forecast -- the $\Fmat$ that the whole risk model runs on.

\begin{keyresult}
The factor covariance forecast is the sample covariance corrected in four ways: recency-weighted (EWMA), serial-correlation-corrected (Newey--West), eigenvalue-de-biased (eigenfactor), and level-corrected (volatility regime). The eigenfactor adjustment is the one that makes the matrix safe to optimise against.
\end{keyresult}

\begin{pitfall}
One honest limitation frames the whole stage: a covariance forecast is still a forecast, and it is calibrated \emph{on average}, not for the next surprise. The pipeline removes the known, structured biases of the sample covariance, but it cannot anticipate a regime the data has not yet shown -- at the very onset of a crisis the level lags until the fast layers catch up, and the correlation structure, deliberately slow, will be a step behind a sudden change in how factors co-move. The model is also calibrated to be mildly conservative on average (bias $\approx 0.95$), a deliberate lean. None of this is a flaw to be fixed so much as the nature of forecasting volatility: the corrections make the matrix unbiased over a long history and well-behaved in an optimiser, which is the most that can be asked of it, but a risk model is a probabilistic statement, not a guarantee about any single month.
\end{pitfall}

\begin{remark}
None of this changes the stage's place in the model: the factor covariance is the matrix at the centre of $\sigma_p^2=\wvec\T\Xmat\Fmat\Xmat\T\wvec+\wvec\T\Del\wvec$, and the four layers exist to make it a trustworthy forecast rather than a backward-looking measurement. The remaining piece of that identity, the specific-risk diagonal $\Del$, is the subject of its own stage.
\end{remark}

\begin{center}\small\itshape
References: Menchero, Orr \& Wang (2011), \textit{The Barra US Equity Model (USE4)} --- Methodology Notes (Sec. 4); Menchero, Wang \& Orr (2011), \textit{Eigen-adjusted covariance matrices}. Measured quantities are from a personal build run.
\end{center}

\end{document}
